% =====================================================================
%   §4. Применения RSA: аутентификация, цифровая подпись, голосование
% =====================================================================
\section{Применения RSA: аутентификация, цифровая подпись, электронное голосование}
\label{sec:rsa-apps}

В предыдущем параграфе мы научились с помощью RSA \emph{передавать секретное сообщение}. На первый взгляд кажется, что этого уже достаточно для всех задач сетевой безопасности. Однако в действительности шифрование решает лишь одну из проблем --- \term{конфиденциальность}. Помимо неё, в современных системах необходимо обеспечивать ещё как минимум три:

\begin{itemize}
  \item \term{Аутентификацию} --- уверенность, что вы общаетесь именно с тем, за кого себя выдаёт собеседник.
  \item \term{Целостность} --- уверенность, что сообщение не было незаметно изменено по дороге.
  \item \term{Неотказуемость} --- невозможность для отправителя впоследствии отказаться от своего сообщения («это не я отправлял»).
\end{itemize}

Удивительно, но все эти задачи можно решить с помощью \emph{той же самой схемы RSA}, если применить её не «прямо», а наоборот --- сначала закрытым ключом, а потом открытым. Об этом и пойдёт речь в этом параграфе.

% ---------------------------------------------------------------------
\subsection{Аутентификация: «докажи, что это ты»}
\label{subsec:auth}

Представим типичную ситуацию: Алиса заходит в свой личный кабинет на сайте банка. Банк хочет убедиться, что это действительно его клиентка, а не злоумышленник, подсмотревший пароль.

Самый простой подход --- запросить пароль. Но у него есть очевидные недостатки: пароль можно подсмотреть, перехватить, угадать, заставить пользователя его выдать. Кроме того, чтобы проверить пароль, банк должен \emph{хранить} его (или его хеш --- см. §\ref{sec:hashing}) у себя на серверах. А значит, при взломе сервера утечёт сразу множество паролей.

Криптография позволяет построить намного более изящную схему, в которой \emph{секрет вообще никогда не передаётся по сети}.

\begin{importantbox}[title={Важно: принцип «доказательства с нулевым разглашением»}]
Сторона $A$ должна доказать стороне $B$, что обладает неким секретом, \emph{не выдавая никакой информации о самом секрете}. Идея, кажущаяся парадоксальной, на практике вполне реализуема и лежит в основе многих криптографических протоколов.
\end{importantbox}

Простейший пример --- \term{протокол challenge-response} («запрос--ответ») на базе RSA. Пусть у Алисы заведена пара ключей: открытый $(e, n)$ опубликован в системе банка, а закрытый $d$ хранится только у неё (например, в специальном USB-токене или в защищённом хранилище смартфона).

\begin{algorithmbox}[title={Алгоритм: аутентификация через RSA}]
\begin{enumerate}
  \item \textbf{Запрос (challenge).} Банк генерирует \emph{случайное} число $r$ ($0 < r < n$), запоминает его и отсылает Алисе.
  \item \textbf{Ответ (response).} Алиса вычисляет $s = r^{d} \bmod n$, используя свой закрытый ключ, и отправляет $s$ обратно.
  \item \textbf{Проверка.} Банк вычисляет $r' = s^{e} \bmod n$ и сравнивает с тем $r$, которое он отсылал. Если $r' = r$, проверка пройдена.
\end{enumerate}
\end{algorithmbox}

\paragraph*{Почему это работает?} В предыдущем параграфе мы доказали, что для любого $r$ выполняется
\[
  (r^{d})^{e} \equiv r \pmod n.
\]
То есть только тот, кто знает $d$, может «правильно» преобразовать $r$ так, чтобы после открытого преобразования $\cdot^{e}$ получилось снова $r$. Любой другой пользователь, не знающий $d$, может лишь \emph{угадывать} ответ --- а вероятность угадать одно нужное число из миллиарда миллиардов исчезающе мала.

\paragraph*{Почему это \emph{безопасно}?}
\begin{itemize}
  \item Алиса \emph{не передаёт} свой закрытый ключ $d$ --- он остаётся у неё.
  \item Каждый раз $r$ \emph{случайное и новое}. Поэтому даже если злоумышленник перехватит весь сеанс и услышит пару $(r, s)$, в следующий раз банк пришлёт другое $r$, и старый ответ $s$ не подойдёт. Атака повторного воспроизведения (replay attack) исключена.
  \item Банк \emph{не хранит секретов клиента}: у него лишь открытый ключ, утечка которого ничем не страшна.
\end{itemize}

\begin{importantbox}[title={Важно: это учебная иллюстрация}]
Описанная схема передаёт \emph{идею} challenge-response, но не является полноценным протоколом с нулевым разглашением. В реальных системах нельзя позволять кому угодно получать $r^{d} \bmod n$ для произвольного $r$: устройство фактически превращается в «оракул» RSA-подписи. Поэтому на практике запрос специально кодируют, а ключи для подписи и аутентификации разделяют.
\end{importantbox}

Схему этого протокола в графическом виде даёт рис.~\ref{fig:p4-1}.

\begin{figure}[H]
\centering
\begin{tikzpicture}[
  >=Stealth,
  every node/.style={font=\small},
  party/.style={draw=prosvBlue, rounded corners, fill=prosvLight, minimum width=2.2cm, minimum height=1.1cm, align=center, thick}
]
  \node[party] (a) at (0,0) {\textbf{Алиса}\\ ключ $d$};
  \node[party] (b) at (8,0) {\textbf{Банк}\\ ключ $e$};

  \draw[->, thick, prosvRed] ([yshift=15pt]b.west) -- node[above, font=\footnotesize]{случайный $r$} ([yshift=15pt]a.east);
  \draw[<-, thick, prosvGreen!70!black] ([yshift=-5pt]b.west) -- node[below, font=\footnotesize]{$s = r^d \bmod n$} ([yshift=-5pt]a.east);

  \node[align=center, font=\footnotesize] at (4, -1.4) {Банк проверяет: $s^e \bmod n \stackrel{?}{=} r$};
\end{tikzpicture}
\caption{Схема аутентификации по запросу--ответу.}
\label{fig:p4-1}
\end{figure}

\begin{interestbox}[title={Это интересно: «запомни всё, чтобы ничего не помнить»}]
В современных смартфонах закрытый ключ $d$ хранится в специальном чипе --- \emph{защищённом анклаве} (Secure Enclave у Apple, Trusted Execution Environment на Android). Этот чип \emph{физически не позволяет считать} закрытый ключ извне. На него можно лишь «передать» число и попросить выполнить с ним операцию $\cdot^d \bmod n$. Так устроены, например, биометрические разблокировки --- отпечаток пальца «разрешает» чипу применить ключ, но не извлекает его.
\end{interestbox}

% ---------------------------------------------------------------------
\subsection{Электронная цифровая подпись}
\label{subsec:eds}

Аутентификация решает задачу «докажи, что ты --- это ты». Теперь рассмотрим родственную, но другую: «докажи, что \emph{именно ты} отправил \emph{именно это} сообщение». Эта задача стоит особенно остро при подписании документов: договоров, заявлений, налоговых деклараций.

Бумажная подпись от руки выполняет сразу две функции: \emph{идентифицирует} автора и \emph{связывает} его с конкретным документом. К сожалению, при простом копировании в электронном виде это свойство теряется: цифровое изображение подписи легко вырезать и приклеить к любому документу. Нужна принципиально иная конструкция.

\begin{definition}{электронная цифровая подпись}{def:p4-1}
\term{Электронная цифровая подпись} (ЭЦП) --- это короткий блок данных $\sigma$, привязываемый к документу $m$, обладающий двумя свойствами:
\begin{itemize}
  \item только владелец секретного ключа может сгенерировать $\sigma$ для произвольного $m$;
  \item любой желающий, зная открытый ключ, может проверить, что $\sigma$ действительно соответствует $m$ и подписан данным владельцем.
\end{itemize}
\end{definition}

Идея конструкции на основе RSA удивительно проста и красива: достаточно поменять местами роли ключей в шифровании.

\begin{algorithmbox}[title={Алгоритм: цифровая подпись RSA}]
Алиса имеет открытый ключ $(e, n)$ и закрытый $d$. Хочет подписать документ $m$ ($0 < m < n$).
\begin{enumerate}
  \item \textbf{Подпись.} Алиса вычисляет
  \[
    \sigma = m^{d} \bmod n
  \]
  и публикует пару $(m, \sigma)$.
  \item \textbf{Проверка.} Любой проверяющий вычисляет $m' = \sigma^{e} \bmod n$ и сравнивает с $m$. Если $m' = m$, подпись подлинна.
\end{enumerate}
\end{algorithmbox}

\paragraph*{Почему это подпись?} По тому же самому соотношению $(m^{d})^{e} \equiv m \pmod n$. Только владелец закрытого ключа $d$ может построить такое $\sigma$, что после возведения в степень $e$ получится исходное $m$. Если бы это умел кто-то другой, он, по сути, умел бы разлагать $n$ на множители --- а это, как мы обсуждали, считается практически невыполнимой задачей.

\paragraph*{Привязка к документу.} В отличие от рукописной подписи, ЭЦП \emph{зависит от содержимого}: если изменить хоть один бит в $m$, изменится и $\sigma$. Невозможно «отделить» подпись от документа и приклеить её к другому: для другого документа потребуется заново вычислять $\cdot^{d} \bmod n$, что без $d$ невыполнимо (рис.~\ref{fig:p4-2}).

\begin{figure}[H]
\centering
\begin{tikzpicture}[
  >=Stealth, every node/.style={font=\small},
  doc/.style={draw=prosvBlue, fill=prosvCream, rounded corners, minimum width=2cm, minimum height=1cm, align=center},
  op/.style={draw=prosvRed, fill=prosvLight, rounded corners, minimum width=2.6cm, minimum height=0.8cm, align=center, thick},
]
  \node[doc] (m1) at (0,0) {Документ $m$};
  \node[op]  (sign) at (3.7,0) {$\sigma = m^d \bmod n$};
  \node[doc] (pair) at (8,0) {$(m,\sigma)$};
  \draw[->, thick] (m1) -- (sign);
  \draw[->, thick] (sign) -- (pair);

  \node[op]  (verify) at (3.7,-2) {$\sigma^e \bmod n$};
  \node[doc] (mres) at (8,-2)   {$m'$};
  \node[align=center, font=\footnotesize] at (8.8,-3) {$m' \stackrel{?}{=} m$};

  \draw[->, thick] (pair.south) -- ++(0,-0.6) -| (verify.north);
  \draw[->, thick] (verify) -- (mres);

  \node[align=left, font=\footnotesize\itshape, prosvBlue] at (-0.5,-2) {Подпись (закр.\ ключ $d$):};
  \node[align=left, font=\footnotesize\itshape, prosvBlue] at (-0.5,-2.6) {Проверка (откр.\ ключ $e$):};
\end{tikzpicture}
\caption{Подписание и проверка электронной цифровой подписи на базе RSA.}
\label{fig:p4-2}
\end{figure}

\paragraph*{Проблема большого документа.} В реальности документ может быть огромным (например, несколько мегабайт). Возводить миллион байтов в степень $d$ по модулю $n$ слишком медленно. Поэтому в практических системах поступают иначе:
\begin{enumerate}
  \item вычисляют \term{хеш} (короткий «отпечаток») документа: $h = H(m)$ --- небольшое число, скажем, длиной в 256 бит;
  \item подписывают именно хеш: $\sigma = h^{d} \bmod n$.
\end{enumerate}

При проверке вычисляют $h' = H(m)$ заново и сравнивают $\sigma^{e} \bmod n$ с $h'$. О хеш-функциях и их математическом устройстве пойдёт речь в §\ref{sec:hashing}.

\begin{importantbox}[title={Важно: это учебная схема}]
Здесь показано лишь математическое ядро RSA-подписи. В реальных стандартах подписывают не «голый» хеш, а специально закодированное значение: применяются схемы дополнения (padding) вроде RSASSA-PSS или RSASSA-PKCS1-v1\_5. Без такого кодирования «наивная» подпись $\sigma = h^{d} \bmod n$ уязвима для ряда атак.
\end{importantbox}

\begin{interestbox}[title={Это интересно: ЭЦП в России}]
В Российской Федерации электронная цифровая подпись юридически приравнена к собственноручной с 2002 года (закон №\,1-ФЗ «Об электронной цифровой подписи», заменён в 2011 году на №\,63-ФЗ «Об электронной подписи»). С помощью ЭЦП сегодня сдают налоговую отчётность, регистрируют недвижимость, участвуют в государственных закупках --- всё это без бумажных носителей. Криптографические алгоритмы, используемые в российских ЭЦП, описаны в стандарте ГОСТ Р 34.10--2012; математически они близки к описанной RSA-схеме, но используют не модулярное возведение в степень, а операции на так называемых \emph{эллиптических кривых}.
\end{interestbox}

% ---------------------------------------------------------------------
\subsection{Электронное голосование}
\label{subsec:voting}

С распространением Интернета возникает естественный вопрос: можно ли перевести голосование --- например, выборы или опросы --- в цифровой формат? На первый взгляд, это просто: создать сайт, попросить каждого выбрать вариант, посчитать голоса. На деле возникает несколько противоречивых требований, и решить их одновременно --- содержательная криптографическая задача.

\begin{importantbox}[title={Важно: требования к электронному голосованию}]
Хорошая система электронного голосования должна одновременно обеспечивать:
\begin{enumerate}
  \item \term{Корректность}: каждый имеющий право проголосовать голосует не более одного раза.
  \item \term{Анонимность}: никто, включая организаторов, не может сопоставить конкретного избирателя с его голосом.
  \item \term{Проверяемость}: избиратель может убедиться, что его голос учтён, а наблюдатели --- что общий подсчёт верен.
  \item \term{Неподдельность}: никто не может вбросить «лишние» голоса.
\end{enumerate}
\end{importantbox}

Заметим противоречие между пунктами 1 и 2: чтобы убедиться, что Алиса голосует не дважды, надо как-то её опознать; но если её опознали, как сохранить анонимность? На бумажных выборах эта задача решается \emph{физически} (список явки и тайная урна разделены), а в цифровом мире --- математически. Ключевой инструмент --- \term{слепая цифровая подпись}, предложенная Дэвидом Чаумом ещё в 1982 году.

\paragraph*{Слепая подпись.} Идея слепой подписи: можно подписать документ, \emph{не зная его содержания}. Этот странный, казалось бы, инструмент имеет естественный аналог в офлайне.

\begin{interestbox}[title={Это интересно: аналогия с конвертом и копиркой}]
Представьте: вы кладёте лист бумаги в конверт, внутри которого вложена копирка. Затем просите начальника поставить подпись \emph{на конверте}. Подпись через копирку переходит и на лист внутри. Раскрыв конверт, вы получаете подписанный документ, причём начальник никогда не видел его содержания.

Слепая RSA-подпись --- это математическая реализация той же идеи: «копирка» --- модулярное умножение на случайный множитель.
\end{interestbox}

Формально схема устроена так. Пусть у центра голосования есть RSA-ключи $(e, n)$ (открытый) и $d$ (закрытый).

\begin{algorithmbox}[title={Алгоритм: слепая подпись RSA}]
Алиса хочет получить подпись центра под её бюллетенем $m$, не раскрывая центру содержимое.
\begin{enumerate}
  \item Алиса выбирает случайное $k$, взаимно простое с $n$, и вычисляет \emph{ослеплённое} сообщение
  \[
    m' = m \cdot k^{e} \bmod n.
  \]
  Отправляет $m'$ центру.

  \item Центр (после проверки, что Алиса действительно имеет право голоса) ставит подпись:
  \[
    \sigma' = (m')^{d} \bmod n = m^{d} \cdot k^{ed} \bmod n = m^{d} \cdot k \bmod n.
  \]
  Возвращает $\sigma'$ Алисе.

  \item Алиса \emph{снимает ослепление}: вычисляет
  \[
    \sigma = \sigma' \cdot k^{-1} \bmod n = m^{d} \bmod n.
  \]
  Получена настоящая RSA-подпись центра под её бюллетенем $m$, хотя центр никогда не видел $m$.
\end{enumerate}
\end{algorithmbox}

\paragraph*{Где здесь теория чисел.} Заметим, что на шаге 3 Алиса делит на $k$, точнее --- умножает на $k^{-1} \bmod n$. Существование \emph{обратного} по модулю $n$ обеспечивается тем, что $\gcd(k, n) = 1$, и его можно эффективно вычислить расширенным алгоритмом Евклида (§\ref{sec:numtheory}). Кроме того, в выкладке использовалось $k^{ed} \equiv k \pmod n$ --- следствие малой теоремы Ферма (точнее, её обобщения --- теоремы Эйлера, гласящей, что $k^{\Eul(n)} \equiv 1 \pmod n$ для $\gcd(k, n) = 1$).

\paragraph*{Сам протокол голосования.} На базе слепой подписи строится следующая схема (в упрощённом виде; подробности --- в книге Музыкантского и Фурина).

\begin{algorithmbox}[title={Алгоритм: электронное голосование}]
\begin{enumerate}
  \item \emph{Подготовка.} Алиса формирует бюллетень $m$ --- например, число, кодирующее её выбор, плюс случайный «уникальный идентификатор» $r$, чтобы исключить совпадение бюллетеней.
  \item \emph{Получение подписи.} Алиса ослепляет $m$, отправляет $m'$ \term{Регистратору}. Регистратор проверяет по своему списку, что Алиса имеет право голосовать и ещё не получала подпись, и возвращает $\sigma'$. Алиса снимает ослепление и получает $(m, \sigma)$.
  \item \emph{Анонимная отправка.} Через анонимный канал (например, через сеть посредников или с другого устройства) Алиса отправляет $(m, \sigma)$ \term{Счётчику}.
  \item \emph{Подсчёт.} Счётчик проверяет подпись $\sigma$: если она верна и идентификатор $r$ ранее не встречался --- голос засчитан.
\end{enumerate}
\end{algorithmbox}

\paragraph*{Почему это работает.}
\begin{itemize}
  \item \emph{Корректность.} Регистратор контролирует, что одна и та же Алиса не получит подпись дважды.
  \item \emph{Анонимность.} Регистратор знает, кто такая Алиса, но не видит содержимого $m$ (оно ослеплено). Счётчик видит $m$ и $\sigma$, но не знает, кому принадлежит подпись --- лишь то, что она поставлена настоящим Регистратором.
  \item \emph{Неподдельность.} Без знания закрытого ключа $d$ невозможно сгенерировать $\sigma$ к произвольному $m$ --- это потребовало бы взлома RSA.
  \item \emph{Проверяемость.} Список $\{(m, \sigma)\}$ публикуется. Алиса может убедиться, что её $(m,\sigma)$ в списке; любой наблюдатель может пересчитать сумму, не зная имён.
\end{itemize}

Общую идею протокола поясняет рис.~\ref{fig:p4-3}.

\begin{figure}[H]
\centering
\begin{tikzpicture}[
  >=Stealth, every node/.style={font=\footnotesize},
  voter/.style={draw=prosvBlue, fill=prosvLight, rounded corners, minimum width=2.4cm, minimum height=1cm, align=center, thick},
  server/.style={draw=prosvRed, fill=prosvCream, rounded corners, minimum width=2.4cm, minimum height=1cm, align=center, thick},
]
  \node[voter] (alice) at (0, 0) {Алиса};
  \node[server] (reg) at (5, 1.5) {Регистратор\\ {\scriptsize ключ $d$}};
  \node[server] (count) at (5, -1.5) {Счётчик};

  \draw[->, thick, prosvRed!80] (alice.north east) -- node[above, sloped]{$m' = m \cdot k^e$} (reg.west);
  \draw[<-, thick, prosvRed!80] ([yshift=-15pt]alice.east) -- node[below, sloped]{$\sigma' = (m')^d$} ([yshift=-15pt]reg.west);
  \draw[->, thick, prosvGreen!60!black] (alice.south east) -- node[below, sloped]{анонимно: $(m, \sigma)$} (count.west);
\end{tikzpicture}
\caption{Электронное голосование на базе слепой подписи. С Регистратором Алиса работает «лицом к лицу» (он знает её, но не видит бюллетень). К Счётчику бюллетень приходит анонимно, но с проверяемой подписью Регистратора.}
\label{fig:p4-3}
\end{figure}

\begin{interestbox}[title={Это интересно: голосование в реальности}]
Системы электронного голосования на основе слепой подписи (и более сложных схем) уже применялись на практике: например, в Эстонии часть выборов проходит через интернет с 2005 года, в России на муниципальных выборах в Москве с 2019 года использовалась система на базе криптографических конструкций (правда, на основе не только RSA, но и других схем --- с гомоморфным шифрованием и блокчейном). Подробное обсуждение разных схем электронного голосования можно найти в книге Музыкантского и Фурина «Лекции по криптографии».
\end{interestbox}

% ---------------------------------------------------------------------
\subsection{Итог параграфа}

Мы увидели, что одна и та же конструкция RSA --- модулярное возведение в степень --- решает четыре совершенно разные задачи:
\[
  \boxed{\;\text{шифрование}\;} \qquad
  \boxed{\;\text{аутентификация}\;} \qquad
  \boxed{\;\text{цифровая подпись}\;} \qquad
  \boxed{\;\text{слепая подпись}\;}
\]
Различие лишь в том, \emph{в каком порядке} применяются операции $\cdot^{e}$ и $\cdot^{d}$ и какие именно числа подставляются на вход. В этом одна из главных красот современной криптографии: глубокая математическая идея даёт целый каскад прикладных решений.

\begin{tasksbox}
\begin{enumerate}
  \item Объясните своими словами, в чём принципиальная разница между «шифрованием с открытым ключом» и «цифровой подписью», если в обеих схемах используются одни и те же ключи и одна и та же формула.
  \item Почему в протоколе аутентификации банк каждый раз присылает \emph{новое случайное} $r$, а не одно и то же?
  \item В схеме ЭЦП Алиса публикует пару $(m, \sigma)$. Что произойдёт, если она опубликует пару $(m', \sigma)$ с изменённым $m' \ne m$? Сможет ли она тем самым обмануть проверяющего?
  \item* В схеме слепой подписи покажите подробной выкладкой, что после снятия ослепления получается именно $m^d \bmod n$ (используйте $k^{ed} \equiv k \pmod n$).
  \item* Предложите простую (возможно, наивную) атаку на схему голосования, если из протокола убрать пункт о случайном идентификаторе $r$ в бюллетене.
  \item* Подумайте, какие ещё свойства бумажного бюллетеня (например, возможность переголосовать, защита от принуждения) электронное голосование решает плохо. Каковы возможные подходы?
\end{enumerate}
\end{tasksbox}
